\documentclass[a4paper, spanish]{article}

\usepackage[top=2cm, bottom=2cm, left=3cm, right=2.5cm]{geometry}

\usepackage{listings}
\usepackage{xcolor}

\definecolor{codegreen}{rgb}{0,0.6,0}
\definecolor{codegray}{rgb}{0.5,0.5,0.5}
\definecolor{codepurple}{rgb}{0.58,0,0.82}
\definecolor{backcolour}{rgb}{0.95,0.95,0.92}

\lstdefinestyle{mystyle}{
    backgroundcolor=\color{backcolour},   
    commentstyle=\color{codegreen},
    keywordstyle=\color{magenta},
    numberstyle=\tiny\color{codegray},
    stringstyle=\color{codepurple},
    basicstyle=\ttfamily\footnotesize,
    breakatwhitespace=false,         
    breaklines=true,                 
    captionpos=b,                    
    keepspaces=true,                 
    numbers=left,                    
    numbersep=5pt,                  
    showspaces=false,                
    showstringspaces=false,
    showtabs=false,                  
    tabsize=2
}


\usepackage[T1]{fontenc}

% \usepackage[utf8]{inputenc} is no longer required (since 2018)

% Spanish-specific commands
\usepackage[spanish]{babel}

% HEADER CONFIGURATION
\usepackage{fancyhdr}

\usepackage{float}

\usepackage[hidelinks]{hyperref}

\usepackage{graphicx}
\graphicspath{ {img/} }

% Adjust the header height to fit your logo without errors
\setlength{\headheight}{32pt} 
\addtolength{\topmargin}{-16pt}

\pagestyle{fancy}
\fancyhf{} % Clear default settings

% Insert logo on the left (adjust height or width to scale)
\fancyhead[L]{\includegraphics[height=18pt]{logo_utp.png}}

% Add text to the right and center if needed
\fancyhead[R]{\nouppercase{\leftmark}}
% \fancyhead[R]{\nouppercase{\rightmark}}

\fancyfoot[R]{\thepage}


\title{Diagnóstico e Implementación de Sistemas de Información Empresarial Para la Empresa MG Solutions S.A.C}
\author{Huanca Fernandez Kevin, Llacsahuanga Buques Diego}

\lstset{style=mystyle}

\begin{document}
\begin{titlepage}

    % 1. LOGO ALINEADO A LA DERECHA
    \begin{flushright}
        % Reemplaza 'logo_utp.png' con el nombre real de tu archivo
        \includegraphics[height=32px]{logo_utp.png} 
    \end{flushright}

    \vspace{3cm} % Espacio entre el logo y los autores

    \begin{flushleft}
    % 2. AUTORES Y TÍTULO ALINEADOS A LA IZQUIERDA
        {\Large \textbf{Diagnóstico e Implementación de Sistemas de Información Empresarial Para la Empresa MG Solutions S.A.C}}\\[1cm]
        
        {\large Huanca Fernández Kevin, Llacsahuanga Buques Diego}\\[0.8cm]
        
        %{\Large Enter your subtitle, comment if not applicable (\% \textbackslash subtitle)}
    \end{flushleft}

    \vfill % Esto empuja automáticamente todo el bloque siguiente hacia el final de la página

    % 3. METADATOS INFERIORES ALINEADOS A LA IZQUIERDA
    \begin{flushleft}
        % \footnotesize % Reduce el tamaño de letra para esta sección
        \textbf{Curso}\\
        Sistemas de Información Empresarial (44576)\\[0.5cm]

        \textbf{Docente}\\
        Nicho Gomez, Miluska\\[0.5cm]

        Sábado 18 de Julio, 2026
    \end{flushleft}
\end{titlepage}

  \section{Implementación del Nodo IoT}
  
  Para la implementación del nodo IoT se utilizó un microcontrolador
  ESP32 que se utilizará para recibir mediciones desde
    la nube y un sensor ultrasónico para medir la próximidad
    de objetos. A continuación se detalla la lista de materiales
    utilizadas para la elaboración.

    \begin{itemize}
      \item \textbf{ESP32 WROOM-32}: es un microcontrolador
      con Wifi y Bluetooth incluído, ideal para recibir lecturas
      de sensores y enviarlas a través de internet.
      \item \textbf{HC-SR04}: sensor ultrasónico utilizable para
      medir la distancia en centímetros utilizando la velocidad del
      sonido con la siguiente fórmula.
      \begin{equation}
        D = \frac{T * 0.034}{2}
      \end{equation}
      \item \textbf{Protoboard}: utilazada para conectar los
      sensores al controlador.
      \item 4 jumpers macho-macho.
      \item Un cable Micro USB a USB.
    \end{itemize}

    Posteriormente programamos el ESP32 con el siguiente código,
    escrito desde Arduino IDE.

  \lstinputlisting[language=C]{src/esp32-to-cloud.ino}

```latex
\section{Configuración del Broker MQTT}

Para gestionar la comunicación entre el nodo IoT y los servicios de almacenamiento, se configuró EMQX 5.8.6 mediante Docker Compose. Se habilitaron los puertos 1883 para la comunicación MQTT y 18083 para acceder al panel de administración. Asimismo, se configuró la autenticación mediante usuario y contraseña para controlar el acceso al broker.

La arquitectura se complementó con InfluxDB para almacenar las mediciones, Node-RED para procesar los mensajes MQTT y Grafana para visualizar los datos. Todos los servicios se conectaron mediante una red Docker compartida. Se adjunta el código de Docker Compose con el despliegue de todos los servicios.

  

\section{Publicación de Datos vía MQTT}

El ESP32 se programó mediante Arduino IDE utilizando las bibliotecas \texttt{WiFi.h} y \texttt{PubSubClient.h}. El dispositivo se conecta a la red Wi-Fi y establece una conexión autenticada con el broker EMQX mediante su dirección IP y el puerto 1883.

El sensor ultrasónico HC-SR04 mide la distancia de los objetos utilizando el tiempo de retorno de la señal sonora. Las mediciones se realizan cada dos segundos y se publican en el tópico \texttt{ta3/distancia} mediante mensajes JSON con el siguiente formato:

\begin{lstlisting}
{"distancia_cm": 25.43}
\end{lstlisting}

La transmisión se verifica mediante el monitor serial del Arduino IDE y un cliente MQTT suscrito al tópico correspondiente, comprobando la recepción continua de los mensajes publicados.

\section{Almacenamiento y Visualización}

Para almacenar las mediciones, se utilizó InfluxDB 2 como base de datos de series temporales. Node-RED recibe los mensajes publicados en el tópico MQTT, extrae el valor de distancia y registra las mediciones en el bucket configurado de InfluxDB, utilizando las credenciales definidas en las variables de entorno de Docker Compose.

Finalmente, se configuró Grafana para consultar los datos almacenados en InfluxDB y representarlos mediante un dashboard. El panel permite visualizar la evolución de la distancia medida a lo largo del tiempo, facilitando el seguimiento de las mediciones y su actualización conforme se reciben nuevos registros.

\end{document}